\chapter{Results}
\label{chap:results}

This chapter is intentionally structured as a \emph{results-writing plan} for the final thesis draft.
Its purpose is to fix the narrative order, the evidence hierarchy, and the claim boundaries before the final quantitative tables and figures are inserted.
The chapter should be written as a progression from baseline operating regimes to the retained controller line, and then to generalization and limitations-aware interpretation.

\section{Chapter purpose and writing order}
\label{sec:results-purpose}

The chapter should answer the following questions in order:

\begin{enumerate}
    \item What does normal operation look like under the retained benchmark setup?
    \item How much does performance degrade under the retained crunch setting?
    \item Within the crunch setting, how does the controller line improve from early to late versions?
    \item Why is \texttt{v6g\_v6d\_compute300} the retained final thesis candidate?
    \item How much of the final result appears to transfer across depots?
    \item What does the retained runtime-policy follow-up show?
    \item What changes when the retained policy is moved to the larger \texttt{DATA003} east/west depots?
    \item What should \emph{not} be claimed from the current evidence?
\end{enumerate}

The final writing order should therefore remain:

\begin{enumerate}
    \item \texttt{EXP00} normal reference
    \item \texttt{EXP01} crunch baseline degradation
    \item \texttt{Scenario1} controller progression:
    \begin{enumerate}
        \item \texttt{v2}
        \item \texttt{v4}
        \item \texttt{v5}
        \item \texttt{v6f}
        \item \texttt{v6g}
    \end{enumerate}
    \item OOD depot validation for the retained final controller
    \item Runtime-policy follow-up
    \item Auxiliary \texttt{DATA003} scaling evidence
    \item Counterexample and non-mainline variants
    \item Synthesis and cautious interpretation
\end{enumerate}

\section{Baseline regime comparison: \texttt{EXP00} vs \texttt{EXP01}}
\label{sec:results-baselines}

\subsection{\texttt{EXP00}: normal-operation reference}
\label{sec:results-exp00}

This subsection should establish the non-stressed reference level of the retained system.
It should describe the expected performance envelope under normal operations and define the baseline interpretation for later comparisons.

\paragraph{What to show.}
Insert the following evidence types:
\begin{itemize}
    \item aggregate table for service rate, failed orders, penalized cost, and runtime across seeds;
    \item short textual summary of the normal operating regime;
    \item one compact figure or table showing day-level stability if this helps anchor later comparison.
\end{itemize}

\paragraph{What to argue.}
The purpose here is not to claim novelty, but to define the operational reference point against which the crunch regime will be judged.
The current aggregate evidence supports a near-perfect normal-operation envelope on the retained Herlev benchmark and very high BAU performance on the larger \texttt{DATA003} depots.
In the current 10-seed aggregates, the retained Herlev \texttt{EXP00} baseline achieves `100\%` service rate, `0.0` failed orders, and penalized cost `6650.63`, while the \texttt{DATA003 west} BAU variants also reach `100\%` service rate and `0` average failed orders.

\paragraph{Writing constraint.}
Do not over-interpret \texttt{EXP00}. It is a reference regime, not a mechanism proof.

\subsection{\texttt{EXP01}: crunch baseline degradation}
\label{sec:results-exp01}

This subsection should quantify the retained pressure setting and explain how the system degrades relative to \texttt{EXP00}.

\paragraph{What to show.}
Insert:
\begin{itemize}
    \item direct side-by-side comparison of \texttt{EXP00} and \texttt{EXP01};
    \item mean differences in service rate, failed orders, and penalized cost;
    \item if useful, a small day-profile figure showing where the degradation becomes visible.
\end{itemize}

\paragraph{What to argue.}
This subsection should establish that \texttt{EXP01} is the retained thesis pressure regime.
It defines the stress context under which the controller line is evaluated.
On the retained Herlev benchmark, the current 10-seed aggregate shows that \texttt{EXP01} baseline performance falls to `96.7816\%` service rate with `33.6` average failed orders and penalized cost `11427.01`, relative to the \texttt{EXP00} reference of `100\%`, `0.0`, and `6650.63`.
This corresponds to a `-3.2184` percentage-point service-rate drop, `+33.6` failed orders, and `+4776.38` penalized cost.

\paragraph{Writing constraint.}
The comparison should be framed as a regime comparison, not as a causal proof of why degradation occurs.

\section{Main method line within \texttt{EXP01 / Scenario1}}
\label{sec:results-scenario1}

This section is the core of the chapter.
It should present the retained controller progression as the main method narrative:
\[
\texttt{v2} \rightarrow \texttt{v4} \rightarrow \texttt{v5} \rightarrow \texttt{v6f} \rightarrow \texttt{v6g}.
\]

The section should use a consistent KPI set throughout:
\begin{itemize}
    \item service rate,
    \item failed orders,
    \item penalized cost,
    \item runtime or compute budget context,
    \item selected day-level or failure-profile diagnostics where needed.
\end{itemize}

\subsection{From \texttt{v2} to \texttt{v4}: commitment logic}
\label{sec:results-v2-v4}

This subsection should explain the first major retained improvement step.

\paragraph{What to show.}
\begin{itemize}
    \item comparison table: \texttt{v2} vs \texttt{v4};
    \item a short interpretation of the shift from earlier robust control to event-driven commitment logic;
    \item one diagnostic if available showing why \texttt{v4} improves the mainline outcome.
\end{itemize}

\paragraph{What to argue.}
The main claim should be that event-driven commitment improves the retained \texttt{Scenario1} line relative to the earlier failure-first baseline.

\paragraph{Writing constraint.}
Do not claim that this step alone solves the stress setting.
It should be presented as an early but important improvement.

\subsection{From \texttt{v4} to \texttt{v5}: risk-budgeted improvement}
\label{sec:results-v4-v5}

This subsection should explain why the thesis line moves from commitment logic to a more structured risk-budgeted controller.

\paragraph{What to show.}
\begin{itemize}
    \item comparison table: \texttt{v4} vs \texttt{v5};
    \item one concise interpretation of how admission / commitment / release logic changes;
    \item optional failure-profile summary if it clarifies the gain.
\end{itemize}

\paragraph{What to argue.}
The safe result interpretation is that \texttt{v5} improves the retained main line and brings the controller closer to the final performance envelope.
In the current aggregate, \texttt{v5\_compute300} reaches `97.3659\%` service rate and `27.5` failed orders, improving on \texttt{v4\_compute300} at `97.1264\%` and `30.0`.

\paragraph{Writing constraint.}
Avoid inflating this step into a final solution.
It is an intermediate retained milestone.

\subsection{From \texttt{v5} to \texttt{v6f}: execution-aware transition}
\label{sec:results-v5-v6f}

This subsection is important because it marks the mechanism transition in the retained storyline.

\paragraph{What to show.}
\begin{itemize}
    \item comparison table: \texttt{v5} vs \texttt{v6f};
    \item evidence that deadline-day drops or execution-related failures become better controlled;
    \item one diagnostic or case interpretation if available.
\end{itemize}

\paragraph{What to argue.}
The main point is not only that \texttt{v6f} improves the KPI line, but that it changes the interpretation of the remaining gap.
At this stage the bottleneck should be written as increasingly execution-aware rather than purely phase/value-selection driven.
The current 10-seed aggregate supports this transition: \texttt{v6f\_compute300} reaches `97.8161\%` service rate with `22.8` failed orders, while the \texttt{v6f\_v6d\_compute300} refit remains close at `97.7969\%` and `23.0`.

\paragraph{Writing constraint.}
Do not over-claim a full causal decomposition.
The safe wording is that \texttt{v6f} is the retained transition step that motivates the final \texttt{v6g} refinement.

\subsection{From \texttt{v6f} to \texttt{v6g}: final retained controller}
\label{sec:results-v6f-v6g}

This subsection is the climax of the results chapter and should make clear why \texttt{v6g} is the retained final thesis candidate.

\paragraph{What to show.}
\begin{itemize}
    \item final ranking table including at least \texttt{v6f} and the retained \texttt{v6g} endpoint;
    \item exact headline KPI values for the retained final candidate;
    \item short explanation of the role of deadline reservation in the final improvement;
    \item optional figure highlighting the remaining gap closure on the main depot.
\end{itemize}

\paragraph{What to argue.}
The final safe claim should identify one canonical main candidate and then separate supplementary follow-ups:
\begin{quote}
Among the retained Herlev mainline controllers, \texttt{v6g\_v6d\_compute300} is the final paper-facing thesis candidate, reaching `97.9885\%` service rate with `21.0` average failed orders.
The \texttt{cap05} rollback and \texttt{automode} results can be reported as supplementary engineering follow-ups, but they should not replace the canonical mainline endpoint in the core ranking.
\end{quote}

\paragraph{Writing constraint.}
Do not generalize this into a universal proof of optimality, physical capacity saturation, or causal completeness.
Keep the statement scoped to the retained experimental setting.

\section{Final ranking and retained endpoint summary}
\label{sec:results-final-ranking}

This section should consolidate the retained mainline comparison into one compact ranking table and one short interpretive paragraph.

\paragraph{Required content.}
Include a table with the retained top endpoints, ideally:
\begin{itemize}
    \item \texttt{v6g\_v6d\_compute300} as the final main candidate,
    \item \texttt{v6f} as the key transition comparator.
    \item \texttt{v5} as the last clearly pre-\texttt{v6f} retained comparator.
\end{itemize}

\paragraph{Interpretive goal.}
This section should let the reader leave the mainline comparison with a single clear conclusion:
the thesis line converges on \texttt{v6g\_v6d\_compute300}, not because every historical variant was inferior in every respect, but because it provides the cleanest retained mechanism closure under the writing-safe scope.

\section{OOD depot validation}
\label{sec:results-ood}

This section should present the cross-depot validation of the retained final controller.

\subsection{Purpose of the OOD study}
\label{sec:results-ood-purpose}

The OOD subsection is not meant to redefine the whole thesis.
Its purpose is to test whether the retained final controller remains strong when the depot changes while the scenario logic is held fixed.

\subsection{What to show}
\label{sec:results-ood-evidence}

Insert:
\begin{itemize}
    \item a compact table for Aalborg, Odense, and Aabyhoj;
    \item service rate and failed orders for the retained final controller;
    \item one sentence reminding the reader that this is an OOD validation under the retained \texttt{Scenario1} framing.
\end{itemize}

\subsection{What to argue}
\label{sec:results-ood-argument}

The safe claim is that the retained \texttt{v6g\_v6d\_compute300} endpoint demonstrates strong transfer under the tested depot shifts.
The current 10-seed aggregates support `100\%` service rate and `0` failed orders on Aalborg, Odense, and Aabyhoj.
The associated penalized costs are `235.24` for Aalborg, `376.69` for Odense, and `399.63` for Aabyhoj.

\subsection{Writing constraint}
\label{sec:results-ood-constraint}

Do not turn this into a universal generalization claim.
It supports robustness and transfer within the tested depots, not a proof of universal deployment behavior.
If the `cap05 / dyn60 / dyn90` OOD sweeps are mentioned, they should be explicitly labeled as `5`-seed supplementary sensitivity checks rather than primary evidence.

\section{Runtime-Policy Follow-Up}
\label{sec:results-runtime-policy}

This section should remain clearly subordinate to the retained Herlev mainline.
Its purpose is not to redefine the final controller, but to show how a practical runtime-policy wrapper behaves when the retained \texttt{v6g\_v6d} controller is already fixed.

\subsection{What to show}
\label{sec:results-runtime-show}

Insert:
\begin{itemize}
    \item a compact comparison between \texttt{v6g\_automode} and \texttt{v6g\_fixed300\_control};
    \item one short explanation of the mode selector: \texttt{dyn60}, \texttt{dyn90}, and \texttt{fixed300};
    \item one sentence clarifying that the current implementation is an engineering runtime-policy layer rather than a new controller family.
\end{itemize}

\subsection{What to argue}
\label{sec:results-runtime-argue}

The safe claim is that \texttt{automode} is the sole retained dynamic-time follow-up.
It can be used to show how runtime is adapted to pressure signals, but it should not displace \texttt{v6g\_v6d\_compute300} as the thesis main candidate.
If numbers are reported, they should be framed as supplementary Herlev evidence: the current 10-seed aggregate for \texttt{automode} reaches `98.1992\%` service rate with `18.8` failed orders, while the \texttt{cap05} fixed-\num{300} rollback reaches `98.3238\%` with `17.5`.

\subsection{Writing constraint}
\label{sec:results-runtime-constraint}

Do not present \texttt{automode} as the final method contribution.
Do not claim that the current runtime selector already constitutes a fully general intelligent compute allocator; in the present implementation it is a retained engineering rule layer driven by pressure signals.

\section{Auxiliary \texttt{DATA003} Scaling Evidence}
\label{sec:results-data003}

This section should remain secondary to the retained Herlev mainline.
Its purpose is to show how the retained policy family behaves on the larger \texttt{DATA003} east/west depots under the same \texttt{EXP00} / \texttt{EXP01} regime logic.

\subsection{What to show}
\label{sec:results-data003-show}

Insert:
\begin{itemize}
    \item one compact BAU table for \texttt{DATA003 east} and \texttt{west};
    \item one compact crunch table for \texttt{DATA003 east} and \texttt{west};
    \item one sentence highlighting that `best service rate` and `lowest penalized cost` coincide on some lines but not on others.
\end{itemize}

\subsection{What to argue}
\label{sec:results-data003-argue}

The safe interpretation is instance-dependent.
For \texttt{DATA003 east}, the retained auxiliary line is the \texttt{east-dyn90} run, which reaches `99.8898\%` service rate, `3.6` failed orders, and penalized cost `59911.87`.
For \texttt{DATA003 west}, the retained auxiliary line is the \texttt{west-w16h} run, which reaches `99.4748\%` service rate, `23.9` failed orders, and penalized cost `181414.89`.
These two lines should be written as controlled auxiliary evidence for larger-depot behavior, not as a second multi-branch ranking problem.

\subsection{Writing constraint}
\label{sec:results-data003-constraint}

Do not rank \texttt{DATA003 east/west} together with the retained Herlev mainline in a single best-model table.
Do not use partially completed endpoints as final evidence; in particular, the current \texttt{w16h\_dyn90} east run and the two west probe directories without \texttt{summary\_final.json} should remain outside paper-facing aggregate claims.

\section{Counterexamples, excluded lines, and non-mainline variants}
\label{sec:results-counterexamples}

This section should remain brief and controlled.
Its purpose is to show disciplined selection of the final thesis line.

\subsection{\texttt{v6b3} as a counterexample}
\label{sec:results-v6b3}

The key safe interpretation is that diversification can improve stress-side robustness while still failing to become the retained best mainline controller.
This subsection should therefore be written explicitly as a counterexample, not as a competing final endpoint.

\subsection{\texttt{v6h} as follow-up, not thesis endpoint}
\label{sec:results-v6h}

If mentioned, \texttt{v6h} should appear only as a narrow solver-side follow-up.
It should not be written as a required part of the final thesis contribution.

\subsection{Historical or removed lines}
\label{sec:results-historical-lines}

Other historical variants should be omitted or summarized in one sentence at most.
The chapter should stay focused on the retained mainline story.

\section{Cross-section synthesis}
\label{sec:results-synthesis}

This section should summarize the chapter in three layers.

\subsection{Operational layer}
\label{sec:results-synthesis-operational}

State how the retained system behaves under:
\begin{itemize}
    \item normal operation,
    \item crunch degradation,
    \item controller improvement under stress.
\end{itemize}

\subsection{Method layer}
\label{sec:results-synthesis-method}

State how the retained method progression should be interpreted:
\begin{itemize}
    \item \texttt{v2} establishes the early robust baseline,
    \item \texttt{v4} improves commitment logic,
    \item \texttt{v5} strengthens risk-budgeted control,
    \item \texttt{v6f} marks the execution-aware transition,
    \item \texttt{v6g} closes the retained mainline and becomes the final candidate.
\end{itemize}

\subsection{Evidence layer}
\label{sec:results-synthesis-evidence}

State explicitly that the chapter supports:
\begin{itemize}
    \item a retained best mainline endpoint,
    \item a strong tested OOD outcome,
    \item a retained runtime-policy follow-up that remains secondary to the main controller line,
    \item auxiliary \texttt{DATA003} scaling evidence with regime-dependent behavior,
    \item a disciplined distinction between mainline and non-mainline variants.
\end{itemize}

\section{Writing guardrails for the final draft}
\label{sec:results-guardrails}

The final version of this chapter must respect the following rules:

\begin{enumerate}
    \item Do not write \texttt{v6b3} as the best model.
    \item Do not write \texttt{v6c} as an active valid route.
    \item Do not write \texttt{v6h} as the final thesis result unless the retained writing map is explicitly updated.
    \item Do not use \texttt{EXP00} and \texttt{EXP01} as causal mechanism proof; use them as retained regime definitions.
    \item Do not write the current model-side execution limit as a proven real-world physical limit.
    \item Do not promote \texttt{cap05} or \texttt{automode} into the main controller ranking.
    \item Do not rank \texttt{DATA003 east/west} and Herlev in the same final-controller table.
    \item Do not use partially completed endpoints as paper-facing final evidence.
    \item Keep the chapter centered on the retained order:
    \[
    \begin{aligned}
    \texttt{EXP00} &\rightarrow \texttt{EXP01} \rightarrow \texttt{v2} \rightarrow \texttt{v4} \rightarrow \texttt{v5} \\
    &\rightarrow \texttt{v6f} \rightarrow \texttt{v6g\_v6d\_compute300} \rightarrow \text{OOD} \\
    &\rightarrow \text{runtime-policy} \rightarrow \texttt{DATA003}.
    \end{aligned}
    \]
\end{enumerate}

\section{Items to insert during final drafting}
\label{sec:results-todo}

Before this chapter is finalized, insert the following:

\begin{itemize}
    \item final baseline table for \texttt{EXP00} and \texttt{EXP01};
    \item controller-line comparison table for \texttt{v2}, \texttt{v4}, \texttt{v5}, \texttt{v6f}, and \texttt{v6g};
    \item final ranking table with the retained top endpoints;
    \item OOD depot table for Aalborg, Odense, and Aabyhoj;
    \item one compact auxiliary table for \texttt{DATA003 east/west} BAU and crunch results;
    \item one compact figure for baseline regime comparison;
    \item one compact figure for controller progression;
    \item optional diagnostic figure for failure-profile or day-level improvement if it helps explain the \texttt{v5} to \texttt{v6f} to \texttt{v6g} transition.
\end{itemize}

\section{Closing paragraph template}
\label{sec:results-closing-template}

The chapter should end with a short synthesis paragraph close to the following logic:

\begin{quote}
The retained evidence shows a clear progression from normal-operation reference to crunch degradation and then to systematic controller improvement within the retained \texttt{Scenario1} line.
Among the documented Herlev endpoints, \texttt{v6g\_v6d\_compute300} is the canonical retained thesis candidate, while \texttt{cap05} and \texttt{automode} remain secondary engineering follow-ups.
The tested OOD results suggest strong transfer across the evaluated depots, and the retained \texttt{DATA003} east/west lines provide controlled auxiliary evidence on larger instances.
At the same time, the evidence should be interpreted conservatively: it supports strong retained endpoints under the modeled settings, not a universal proof of physical-optimal operation.
\end{quote}
